% problem_id: S001-lemma-set-theoretic-pre-equivalence-geometric
% source_id: S001 (Stacks Project)
% source_file: groupoids-quotients.tex
% statement_locator: groupoids-quotients.tex lines 873--890
% proof_locator: groupoids-quotients.tex lines 892--1003
% env: lemma
% id_note: label 在本来源内唯一
% pairing: tight (verified)
% status: complete (statement + author proof verbatim + reason + locators)
%
\begin{lemma}
\label{lemma-set-theoretic-pre-equivalence-geometric}
In the situation of Definition \ref{definition-set-theoretic-equivalence}.
The following are equivalent:
\begin{enumerate}
\item The morphism $j$ is a set-theoretic pre-equivalence relation.
\item The subset $j(|R|) \subset |U \times_B U|$ contains the image of
$|j'|$ for any of the morphisms $j'$ as in Equation (\ref{equation-list}).
\item For every algebraically closed field $k$ over $B$ of sufficiently large
cardinality the subset $j(R(k)) \subset U(k) \times U(k)$ is an equivalence
relation.
\end{enumerate}
If $s, t$ are locally of finite type these are also equivalent to
\begin{enumerate}
\item[(4)] For every algebraically closed field $k$ over $B$
the subset $j(R(k)) \subset U(k) \times U(k)$ is an equivalence relation.
\end{enumerate}
\end{lemma}

\begin{proof}
Assume (2). Let $k$ be an algebraically closed field over $B$.
We are going to show that $\sim_R$ is an equivalence relation.
Suppose that $\overline{u}_i : \Spec(k) \to U$, $i = 1, 2$
are $k$-valued points of $U$. Suppose that $(\overline{u}_1, \overline{u}_2)$
is the image of a $K$-valued point $r \in R(K)$. Consider the
solid commutative diagram
$$
\xymatrix{
\Spec(K') \ar@{..>}[r] \ar@{..>}[d]
&
\Spec(k) \ar[d]_{(\overline{u}_2, \overline{u}_1)} &
\Spec(K) \ar[d] \ar[l] \\
R \ar[r]^-j &
U \times_B U &
R \ar[l]_-{j_{flip}}
}
$$
We also denote $r \in |R|$ the image of $r$.
By assumption the image of $|j_{flip}|$ is contained in the image of
$|j|$, in other words there exists a $r' \in |R|$ such that
$|j|(r') = |j_{flip}|(r)$. But note that $(\overline{u}_2, \overline{u}_1)$
is in the equivalence class that defines $|j|(r')$ (by the commutativity
of the solid part of the diagram). This means there exists a field
extension $K'/k$ and a morphism $r' : \Spec(K) \to R$
(abusively denoted $r'$ as well) with
$j \circ r' = (\overline{u}_2, \overline{u}_1) \circ i$
where $i : \Spec(K') \to \Spec(K)$ is the obvious map.
In other words the dotted part of the diagram commutes.
This proves that $\sim_R$ is a symmetric relation on $U(k)$.
In the similar way, using that the image of $|j_{diag}|$ is contained
in the image of $|j|$ we see that $\sim_R$ is reflexive (details omitted).

\medskip\noindent
To show that $\sim_R$ is transitive assume given
$\overline{u}_i : \Spec(k) \to U$, $i = 1, 2, 3$
and field extensions $K_i/k$ and points
$r_i : \Spec(K_i) \to R$, $i = 1, 2$ such that
$j(r_1) = (\overline{u}_1, \overline{u}_2)$ and
$j(r_1) = (\overline{u}_2, \overline{u}_3)$. Then we may choose a
commutative diagram of fields
$$
\xymatrix{
K & K_2 \ar[l] \\
K_1 \ar[u] & k \ar[l] \ar[u]
}
$$
and we may think of $r_1, r_2 \in R(K)$. We consider the
commutative solid diagram
$$
\xymatrix{
\Spec(K') \ar@{..>}[r] \ar@{..>}[d]
&
\Spec(k) \ar[d]_{(\overline{u}_1, \overline{u}_3)} &
\Spec(K) \ar[d]^{(r_1, r_2)} \ar[l]
\\
R \ar[r]^-j &
U \times_B U &
R \times_{s, U, t} R \ar[l]_-{j_{comp}}
}
$$
By exactly the same reasoning as in the first part of the proof, but
this time using that $|j_{comp}|((r_1, r_2))$ is in the image of $|j|$,
we conclude that a field $K'$ and dotted arrows exist making the
diagram commute. This proves that $\sim_R$ is transitive and concludes
the proof that (2) implies (1).

\medskip\noindent
Assume (1) and let $k$ be an algebraically closed field over $B$ whose
cardinality is larger than $\lambda(R)$, see
Morphisms of Spaces, Lemma \ref{spaces-morphisms-lemma-large-enough}.
Suppose that $\overline{u} \sim_R \overline{u}'$ with
$\overline{u}, \overline{u}' \in U(k)$. By assumption there exists
a point in $|R|$ mapping to $(\overline{u}, \overline{u}') \in |U \times_B U|$.
Hence by
Morphisms of Spaces, Lemma \ref{spaces-morphisms-lemma-large-enough}
we conclude there exists an $\overline{r} \in R(k)$ with
$j(\overline{r}) = (\overline{u}, \overline{u}')$. In this way we see
that (1) implies (3).

\medskip\noindent
Assume (3). Let us show that
$\Im(|j_{comp}|) \subset \Im(|j|)$. Pick any point
$c \in |R \times_{s, U, t} R|$. We may represent this by a morphism
$\overline{c} : \Spec(k) \to R \times_{s, U, t} R$, with $k$ over $B$
having sufficiently large cardinality. By assumption we see that
$j_{comp}(\overline{c}) \in U(k) \times U(k) = (U \times_B U)(k)$
is also the image $j(\overline{r})$ for some $\overline{r} \in R(k)$.
Hence $j_{comp}(c) = j(r)$ in $|U \times_B U|$ as desired (with
$r \in |R|$ the equivalence class of $\overline{r}$). The same argument
shows also that $\Im(|j_{diag}|) \subset \Im(|j|)$ and
$\Im(|j_{flip}|) \subset \Im(|j|)$ (details omitted).
In this way we see that (3) implies (2). At this point we have
shown that (1), (2) and (3) are all equivalent.

\medskip\noindent
It is clear that (4) implies (3) (without any assumptions on $s$, $t$).
To finish the proof of the lemma we show that (1) implies (4) if $s, t$
are locally of finite type. Namely, let $k$ be an algebraically closed
field over $B$. Suppose that $\overline{u} \sim_R \overline{u}'$ with
$\overline{u}, \overline{u}' \in U(k)$. By assumption the algebraic space
$Z = R \times_{j, U \times_B U, (\overline{u}, \overline{u}')} \Spec(k)$
is nonempty. On the other hand, since $j = (t, s)$ is locally of finite type
the morphism $Z \to \Spec(k)$ is locally of finite type as well (use
Morphisms of Spaces, Lemmas \ref{spaces-morphisms-lemma-permanence-finite-type}
and \ref{spaces-morphisms-lemma-base-change-finite-type}).
Hence $Z$ has a $k$ point by
Morphisms of Spaces, Lemma
\ref{spaces-morphisms-lemma-locally-finite-type-surjective-geometric-points}
and we conclude that $(\overline{u}, \overline{u}') \in j(R(k))$
as desired. This finishes the proof of the lemma.
\end{proof}
